\section{\centering Encodings \& Constructions in the \lCalc}

This section presents a number of conventional constructions that represent key ideas common to computer science and other fields. We are specifically referring to natural numbers, branching statements, booleans and lists among others, as well as the operations over them. These constructions can be used to prove relevant statements concerning the expressiveness of the $\la$-calculus.

\begin{note}
  We use our Haskell program to bring some computer science constructions to life, which means that some of the details in this section are implementation dependent. In addition, we use the following type aliases to simplify notation: \texttt{v} \; \texttt{=} \; \texttt{Var;} \texttt{l} \; \texttt{=} \; \texttt{Abs;} \texttt{a} \; \texttt{=} \; \texttt{App}.
\end{note}

\subsubsection{\centering Natural Numbers}

\begin{definition} Peano Axioms, the set of natural numbers $\mathbb{N}$
\begin{itemize}
\item \textbf{(Zero)} \( 0 \in \NAT \)
\item \textbf{(Successor)} \( \forall n \in \NAT, S(n) \in \NAT \)
\item (Initial) \( \forall n \in \NAT, S(n) \neq 0 \) \label{def:peano-nat-initial}
\item (Injectivity) \( \forall n, m \in \NAT, S(n) = S(m) \Rightarrow n = m \)
\item (Induction) \( P(0) \land \forall ( P(n) \Rightarrow P(S(n))) \Rightarrow \forall n,  P (n) \)
\end{itemize}
The first four axioms can be summarized using a graph:
\end{definition}
\begin{example}
  Say for instance we mean to represent $ 3 $ using Peano notation, using the graph, we would come up with $ s(s(s(0))) $, since it is $3$ steps away from $0$.
\end{example}
\begin{definition} To add and multiply natural numbers we use following the rules:
  \[
    n + m = 
    \begin{cases}
      n & \text{if } m = 0 \\
      s(n + m') & \text{if } m = s(m')
    \end{cases}
    \quad
    a \cdot b =
    \begin{cases}
      0 & \text{if } b = 0 \\
      a \cdot b' + a & \text{if } b = s(b')
    \end{cases}
  \]
\end{definition}
\begin{definition} Representation of the set of all naturals in the \lcalc, usually referred to as Church numerals:
  \label{def:church-naturals}
  \begin{align*}
    \overline{\texttt{0}} &\equiv \la f.\la x. x \\
    \texttt{SUCC} &\equiv \la n. \la f. \la x . f ( n f x )
  \end{align*}
  Using our Haskell program, $\overline{\texttt{0}}$ is represented like follows:
  \begin{minted}{Haskell}
    zero  = l "f" (l "x" (v "x"))
  \end{minted}
\end{definition}
\begin{example}
  The church numerals have been defined similarly to the peano axioms, that is, by bringing into existance an initial element and using the succesor function to define the rest of them. Let us apply \texttt{SUCC} a few times to illustrate:
  \begin{align*}
    \overline{\texttt{0}} &\equiv \la f.\la x. x \\
    \overline{\texttt{1}} &\equiv \la f.\la x. f x \\
    \overline{\texttt{2}} &\equiv \la f.\la x. f (f x) \\
    \overline{\texttt{3}} &\equiv \la f.\la x. f (f (f x)) \\
    \overline{\texttt{n}} &\equiv \la f.\la x. f ^ n x
  \end{align*}
  where $ n $ denotes the number of times we append $ f $.
\end{example}
\begin{notation}
  We use $\overline{\text{number}}$ to denote a number is church encoded.
\end{notation}
In our Haskell program we use \texttt{nat} to convert from Haskell's prelude type \texttt{Natural} to type our custom type \texttt{Term}.
\begin{minted}{haskell}
  nat :: Natural :: Term
  nat n = normalize (nat_ n)
      where
          nat_ 0 = zero
          nat_ n = a succ_ (nat_ (n - 1))
\end{minted}
\begin{remark}
  Since \lterms \ are first-class citizens, and $f$ and $x$ are abstracted away, the reality is that numbers are a quite general construct that has many distinct applications. Recall Example \ref{ex:high-order-by-example} where $A \equiv_\alpha \overline{\texttt{2}}$.
\end{remark}
\subsubsection{\centering Basic Arithmetic}
\begin{definition} Now that we have a construction for the naturals, we can define some arithmetic operations:
\begin{align*}
  \texttt{ADD} &\equiv \la m.\la n. \la f. \la x. m f ( n f x) \\
  \texttt{MUL} &\equiv \la f.\la x. \la f. \la x. m (n f) x \\
  \texttt{ISZERO} &\equiv \lambda n. n\,(\lambda x.\texttt{FALSE})\,\texttt{TRUE}\\
  \texttt{PRED} &\equiv \la n.\,\la f.\,\la x.\, \bigl(((n \; (\la g.\,\la h.\, h\,(g\,f))) \; (\la u.\,x)) \; (\la u.\,u))\bigr)
\end{align*}
\texttt{PRED} computes the predecessor of a church encoded natural. $\texttt{PRED} \, \overline{0} = \overline{0}$. We omit further Haskell implementation of constructions from now on since they are just literal translations to our programs representation of \lterms.
\end{definition}
\begin{example} Now we illustrate how addition works by performing $2 + 3$ with the system we defined.
    \[
    \begin{aligned}
      \texttt{ADD}\;\texttt{2}\;\texttt{3} &\equiv (\la m.\la n.\la f.\la x. m\,f\,(n\,f\,x))\;(\la f.\la x. f\,f\,x)\;(\la f.\la x. f\,f\,f\,x) \\
                                           &\curly^* \la f.\la x.\; (\la f.\la x. f\,f\,x)\;f\;((\la f.\la x. f\,f\,f\,x)\;f\;x) \\
                                           &\curly^* \la f.\la x.\; (\la x. f\,f\,x)\;((\la x. f\,f\,f\,x)\;x) \\
                                           &\curly (\la f.\la x.\; f\,f\,f\,f\,f\,x) \\
                                           & \equiv \texttt{5}
    \end{aligned}
  \]
  In Haskell:
  \begin{minted}{shell-session}
  ghci> normalize (a (a mul (nat 2)) (nat 3))
\end{minted}
This unequivocally yields \texttt{(\lambda f.(\lambda x.f(f(f(f(f(fx)))))))}.
\end{example}
\begin{remark}
  Multiplication works similarly to addition. As for \texttt{ISZERO}, the intuition behind it is that if the number is zero, then \(\lambda x.\, \texttt{FALSE}\) is dropped, and the result is \texttt{TRUE}; otherwise, it is applied at least once, which immediately yields \texttt{FALSE}.
  \end{remark}
\begin{remark}
  This proves that the \lcalc \ is capable of arithmetic.
\end{remark}
\subsubsection{\centering Booleans \& Branching}
The next constructions in line are the truth values of boolean logic, ``\textit{true}'' and ``\textit{false}'' for short.
\begin{definition} Boolean truth values:
\(
    \texttt{TRUE} \equiv \lambda t.\lambda f. t \text{ and }
    \texttt{FALSE} \equiv \lambda t.\lambda f. f
\)
\end{definition}
\begin{remark}
  \label{rem:true-false}
  Note how $\TRUE \ M \ N \curly M $ and $ \FALSE \ M \ N \curly N $
\end{remark}
\begin{definition} Conditional, in other words \textit{``if statement''}.
  \(
    \texttt{IF} \equiv \lambda b.\lambda t.\lambda e. b\,t\,e \)
\end{definition}
\begin{note}
  \IF \ statement is just a combinator that happens to behave similarly to a branching statement whenever we pass a church encoded boolean as a first argument.
  \end{note}
\begin{example} On the correct behabiour of $\IF$ i.e. check that it implements branching. $ M, N \in \La $ and $ \IF, \TRUE, \FALSE$ are as usual:
  \[
    \begin{aligned}
      & \IF \ \TRUE \ M \ N \\
      \equiv \ & (\la b.\la t.\la e.\, b\, t\, e) \ \TRUE \ M \ N  \\
      \curly \ & (\la t.\la e.\, \TRUE\ t\, e) \ M \ N  \\
      \curly \ & (\la e.\, \TRUE\ M\, e) \ N  \\
      \curly \ & \TRUE\ M\, N \curly^* M 
    \end{aligned}
    \quad
    \begin{aligned}
      & \IF \ \FALSE \ M \ N \\
      \curly^* \ & \FALSE\ M\, N \curly^* N
    \end{aligned}
  \]
  So the \IF \ just combines $b, t, e$ in such a way that whenever b is a church encoded boolean either $t$ or $e$ is dropped depending on the truth value of $b$. See Remark \ref{rem:true-false}. From this example we extract that $ \la x . x $ is equivalent to \IF \ it terms of behaviour, and thus, we could use $ \IF \equiv \la x . x $ interchangeably.
\end{example}
\begin{remark}
  Using the \IF \ we can implement an universal set of logic gates:
  \begin{align*}
    \NOT &\equiv \IF \, b \, \FALSE \, \TRUE \\
    \AND &\equiv \IF \, a \, (\IF \, b \, \TRUE \, \FALSE) \, \FALSE
  \end{align*}
  This allows us to simulate a nand gate, known to be universal. Thus, the lambda calculus is at least equivalent to propositional logic.
\end{remark}

\subsubsection{\centering Pairs}

A fundamental construct for building data structures is the notion of pairs. 
They allow us to bind two values together and query them independently.

\begin{definition}
  \( \PAIR \equiv \la x.\la y.\la f. f\,x\,y \)
\end{definition}

\begin{remark}
  This definition is convenient since 
  \[
    \PAIR \, M \, N \curly^* \la f . f \, M \, N
  \] 
  which means we can retrieve the first component via 
  \( (\la f . f \, M \, N) \, \TRUE \) 
  and the second via 
  \( (\la f . f \, M \, N) \, \FALSE \). 
  This motivates the following projection operators:
\end{remark}

\begin{definition}
  A \PAIR \ is nothing but an ordered relation $\langle A, B \rangle \in A \times B$.  
  The projections $\pi_1, \pi_2$ are defined as:
  \[
    \begin{aligned}
      \FST &\equiv \la p. p\,(\la x.\la y. x) \\
      \SND &\equiv \la p. p\,(\la x.\la y. y)
    \end{aligned}
  \]
\end{definition}

\begin{proposition}
  A key property of $ \PAIR $ is:
  \[
    \PAIR \, A \, B \equiv_\beta 
    \PAIR \, (\FST \, (\PAIR \, A \, B)) \, (\SND \, (\PAIR \, A \, B))
  \]
\end{proposition}

\subsubsection{\centering Lists}
We move on to how we can now define linear lists.

\begin{definition}
  We define the empty list in the same spirit as \texttt{0} and \FALSE:
  \[
    \NIL \equiv \la f.\la z. z
  \]
  It represents the terminating element of the inductive definition of lists. To build non-empty lists we use the constructor:
  \[
    \texttt{CONS} \equiv \la h.\la t.\la f.\la z. f\,h\,(t\,f\,z)
  \]
  which binds a new element $h$ to a pre-existing list $t$ or to \NIL.
\end{definition}
\begin{definition}
  Some other relevant list operators are:
  \begin{align*}
    \texttt{ISNIL} &\equiv \la l. l\,(\la h.\la t. \FALSE)\,\TRUE \\
    \texttt{HEAD} &\equiv \la l.\,((l \; (\la h.\, \la t.\, h))\; \texttt{NIL}) \\
    \texttt{TAIL} &\equiv \la l.\la c.\la n.l\ (\la h.\la r.\la g.g\ h\ (r\ c))\ (\la c.n)\ (\la h.\la t.t)
  \end{align*}
\end{definition}



